\documentclass[11pt]{article}
\usepackage{fontspec}
\setmainfont{TeX Gyre Pagella}
\usepackage{amsmath}
\usepackage{amssymb}
\usepackage[margin=1.25in]{geometry}
\usepackage{titling}
\usepackage{hyperref}
\hypersetup{colorlinks=true, linkcolor=black, citecolor=black, urlcolor=black}

\title{The Latency of Evidence \\ \large Why Facts Can Exist Before a System Becomes Capable of Knowing Them}
\author{Flyxion \\ \small Independent Researcher}
\date{2026}

\begin{document}
\maketitle

\begin{abstract}
Traces can be preserved long before an epistemic system possesses the concepts, methods, access, or incentives required to use them as evidence for a particular proposition. A record may already sit in a log, a photograph, a measurement, a testimony, or a physical residue, and remain practically unusable not because it is hidden but because nothing yet connects it to a question. This essay develops that claim precisely rather than loosely. It distinguishes a trace's material availability from its recoverability, its recoverability from its warranted evidential bearing on a proposition, its warranted bearing from an institution's willingness to admit it toward a judgment, and its admission from any actual downstream consequence. These four statuses, availability, legibility, admissibility, and consequence, are not mutually independent; each later status ordinarily presupposes the one before it. What the essay argues is narrower and more defensible: earlier statuses do not guarantee later ones, and the gap between them can be large, can persist indefinitely, and can close only when a later interpreter supplies a query, a method, or an institutional willingness that did not previously exist. The essay develops a formal apparatus for this structure, applies it to scientific anomalies, software debugging, forensic and legal evidence, and machine learning, connects it to a companion analysis of context reification in generative systems by treating both as opposite errors of a single evidential support relation, and closes with a qualified argument about preservation: uncertain present relevance does not imply zero expected future value, though this consideration must be weighed against the genuine costs, including privacy, security, storage, and the paradoxical risk that indiscriminate preservation buries the traces that matter under those that do not.
\end{abstract}

\section{The Interval Between Existence and Knowledge}

It is tempting to treat discovery as a moment at which something new comes into being. A previously unnoticed entry in a maintenance log is read for the first time and a mechanical failure is explained; a photograph taken decades earlier is re-examined and a disputed sequence of events is settled; a sample preserved for one purpose is reanalyzed under a later technique and yields a result no one had thought to ask for. In each case it is natural to say that something was found. It is more accurate, though less natural, to say that nothing was found in the sense of having newly come into existence. The log entry, the photograph, and the sample were already there. What changed was the relation between that material and an inquiring system, not the material itself.

This essay treats that relation as having a temporal structure of its own, distinct from the temporal structure of the event the material bears on. A trace may become materially available at one time, recoverable at a later time, warrantedly informative about some proposition at a later time still, and institutionally acted upon later than that, or it may never complete this sequence at all. The remainder of the essay develops this structure carefully enough to say something more precise than that knowledge lags behind evidence. The central correction the essay makes to a cruder version of this idea, developed in response to close criticism of an earlier formulation, is that these stages are not logically independent way-stations a trace visits in arbitrary order. They form a dependency chain in which each later stage ordinarily presupposes the one before it, even though none of the earlier stages guarantees that the later ones will ever be reached.

\end{document}
